\section{Instrument validation}\label{sec:validation}

A simulator paper that reports a stylized-facts battery it passes has, in this literature, reported very little. \Citet{cont2001empirical} says so in the founding paper: the stylized facts are ``usually formulated in terms of qualitative properties of asset returns and may not be precise enough to distinguish among different parametric models.'' The objection has been restated at every subsequent methodological turn. \Citet{fagiolo2019validation} grade empirical validity on four levels, from a Level 0 caricature through qualitative and then quantitative agreement with macro structures to Level 3 quantitative agreement with micro structures, observe that publishing standards now require at least Level 1, and survey the techniques built to ``better discriminate among different ABMs reproducing the same set of stylized facts.'' \Citet{lamperti2018information} is blunter, that ``moments are poor representations of time series' behaviour'': two series can agree on every stylized fact and every moment while their dynamics diverge point for point. Twenty-four years after Cont, LOB-Bench restates the objection for machine learning, that prior work ``relied on qualitative comparisons of stylized facts, making rigorous model comparisons infeasible and hindering research progress'' \citep{nagy2025lobbench}. The standard failure mode of this genre is therefore not a simulator that fails its realism check. It is a simulator that passes one and is degenerate anyway, because the check cannot tell models apart.

Read against that, the headline result of this section is evidence about the instrument rather than a shortfall in it. A public, committed, deterministic generator fails a calibrated three-check conjunction on 23 of 24 seeded panels, against an IID null fixed in source and tuned to no market panel, and the failure is reproducible from a command in \cref{app:commands} by anyone who disbelieves it. A gate a generator can fail is a gate with discriminating power, and in this literature that is rarer than a generator that passes an uncalibrated one. We therefore report the result the way \citet{vyetrenko2020getreal} and \citet{nagy2025lobbench} report theirs: we name the fact that fails (absolute-return autocorrelation, on all three tiers), name the configurations that fail it (the canonical generator everywhere, and the Calm tier on excess kurtosis besides), report the opt-in driver's relative improvement without promoting it to a realism claim, and locate the contribution in the measuring instrument rather than in the thing measured. \Citet{vyetrenko2020getreal} set the register in one sentence: ``there are still significant discrepancies between simulated markets and real ones. However, this work serves as a benchmark against which we can measure future improvement.'' \Citet{platt2018agent} reach the same conclusion from the other direction, that ``traditional stylized fact-centric validation seems unable to detect these potential problems, suggesting that such methods of validation are simply not sufficient.'' On the Fagiolo ladder the canonical tape sits at Level 0 on the axis this section gates and makes no claim above it.

Validation runs on four axes, and the section is organized by them, because the environment fails the first and passes structural checks on the other three. A single verdict about ``realism'' would misreport all four. \Cref{sec:f4} is \emph{unconditional} realism, whether the marginal and autocorrelation structure of the tape resembles the classical catalogue. \Cref{sec:f5} is \emph{response} realism, whether the price responds to flow the way a no-manipulation model requires. \Cref{sec:f6} is \emph{mechanism} realism, whether informed flow costs a maker what the microstructure literature says it should. \Cref{sec:predictability} is \emph{adversarial} realism, whether a public deterministic generator withstands an adversary who attacks the generator rather than the interface. Within each axis, every follow-up unit sits directly after the finding it extends.

F4 grades the canonical position-trading generator; F5 and F6 run on the market-making and market-clearing models, which are separate simulators with their own structural and positive-control checks and are not submitted to F4, as \cref{sec:limitations} states. Dispersion follows three conventions: t-based 95\% intervals over committed per-seed or per-episode vectors for F5 and F6, count-based pass rates for F4, and across-seed standard deviations for the predictability probe, each stated where it is used. Every entry point, seed set and quantity is named where it is reported, the numbers come from JSON and figures committed with their producer, and the commands are in \cref{app:commands}.

\subsection{Unconditional realism: stylized-facts certification (F4)}\label{sec:f4}

For each tier we grade price panels from eight seeded episodes with \texttt{certify\_realism} against three directional bounds from the classical catalogue: positive excess kurtosis, absolute-return autocorrelation above a fixed 95th-percentile IID null for this exact $120\times4$, ten-lag estimator, and aggregation moving \emph{absolute} excess kurtosis toward zero \citep{cont2001empirical}. The IID null is a deterministic 20{,}000-panel simulation fixed in source and tuned to no market panel. Gain/loss skew \citep{cont2001empirical}, timescale asymmetry \citep{zumbach2009time}, and a finite-panel-calibrated Fano exceedance-count ratio are informational. The Fano ratio conditions on the observed exceedance count and divides the Bessel-corrected sample Fano by its exact equal-block IID expectation; it remains ungated because 12 count windows make it too noisy for certification.

\begin{table}[t]
\centering
\caption{Stylized-facts results over 8 seeds. Fano is a conditional-IID-calibrated exploratory ratio and has no pass column; the final rate is the conjunction of the three calibrated directional checks, including aggregation toward zero absolute excess kurtosis. Aggregation passes 8/8 seeds in Calm, 8/8 in Hard, and 8/8 in Extreme.}
\label{tab:f4}
\small
\setlength{\tabcolsep}{5pt}
\begin{tabular}{l rr rr r r}
\toprule
 & \multicolumn{2}{c}{Excess kurtosis} & \multicolumn{2}{c}{$|r|$ autocorr} & {Fano proxy} & {Gated facts} \\
\cmidrule(lr){2-3}\cmidrule(lr){4-5}
Tier & mean & pass & mean & pass & mean & pass rate \\
\midrule
Calm & $-0.97$ & 0/8 & $+0.0002$ & 1/8 & 1.32 & 0.000 \\
Hard & $+16.30$ & 8/8 & $-0.0107$ & 0/8 & 0.93 & 0.000 \\
Extreme & $+11.38$ & 8/8 & $-0.0118$ & 1/8 & 0.91 & 0.125 \\
\bottomrule
\end{tabular}
\end{table}

\Cref{tab:f4} and \cref{fig:f4} give a mixed verdict. Calm is platykurtic (mean excess kurtosis $-0.97$), although its returns move closer to Gaussian in absolute-kurtosis distance on aggregation (mean $+0.57$), and fails the conjunction on all eight seeds. Hard and Extreme pass excess kurtosis and aggregation on all eight seeds, but neither clears the calibrated ACF threshold consistently: Hard passes 0 of 8 and Extreme 1 of 8. The canonical generator therefore passes 1 of 24 seeded panels, and the single binding fact across all three tiers is volatility clustering as the ACF measures it. The Fano ratios are observations and do not decide certification.

The opt-in volatility-clustering driver (\texttt{vol\_clustering}) adds deterministic EMA persistence to the same seeded panel while preserving the zero-knob golden path. Under the calibrated ACF rule at strength $0.5$, the three-check pass rates rise to 0/8, 8/8 and 4/8 for Calm, Hard and Extreme, so the driver moves Hard from failing to passing and Extreme halfway. That is a relative improvement in a named opt-in configuration, and we state it as one: the certified path is the zero-knob canonical generator, and it still fails. The opt-in jump-burst knobs and the exploratory Fano ratio support no remediation claim.

\begin{figure}[t]
\centering
\includegraphics[width=\linewidth]{f4-realism}
\caption{Per-seed stylized-fact values by tier. Only kurtosis, absolute-return autocorrelation and aggregational Gaussianity are gated; Fano is exploratory.}
\label{fig:f4}
\end{figure}

\subsubsection{Calibrating Calm with the existing knobs: a negative result}\label{sec:calm-tails}

The obvious next question is whether the shipped knobs can be tuned into a Calm tier that passes, so we ask it under a rule fixed before the sweep. We sweep the existing opt-in families, \texttt{vol\_clustering} in $\{0,0.3,0.5\}$ and the jump-burst probability, persistence and size grid, for 99 configurations. The protocol is fixed in the producer: a setting must pass the three-check conjunction on at least 7 of 8 diagnostic seeds, have mean realized-volatility ratio at most 1.25 versus default Calm, and pass at least 7 of 8 disjoint confirmation seeds (100--107); among survivors it would minimize normalized RMS return perturbation. The ACF check uses a one-sided 5\% IID-null test for the exact $120\times4$, ten-lag estimator rather than the raw condition ``ACF $>0$''. No cell satisfies all three predeclared steps, so there is no certified Calm preset and no report-band estimate to promote.

The closest low-volatility candidate is \texttt{vol\_clustering}$=0.5$, burst probability 0.02, persistence 0 and size 0.015. It passes 7/8 diagnostic panels but only 5/8 confirmation panels at a mean volatility ratio 1.150. Several settings reach 7/8 on both bands only by exceeding the 1.25 volatility constraint; the strongest near candidate (clustering 0.5, probability 0.02, persistence 0.5, size 0.015) has 8/8 and 7/8 passes but ratio 1.289. We retain the former as a reproducible \texttt{calm\_calibration\_candidate} in the Rust and Python surfaces, pinned by a golden fingerprint, precisely so a future generator revision can be compared against it; it is not a certified configuration (\cref{fig:calm-calibration}).

\begin{figure}[t]
\centering
\includegraphics[width=\linewidth]{f4-calm-calibration}
\caption{Calm calibration sweep. Each point is one of 99 configurations; the dashed line
is 7 of 8 diagnostic passes and hollow points exceed the 1.25 volatility constraint. No
configuration meets the full diagnostic--confirmation--volatility rule under the
finite-panel-calibrated ACF gate.}
\label{fig:calm-calibration}
\end{figure}

\subsection{Response realism: the manipulation boundary and size response (F5)}\label{sec:f5}

Everything in this subsection is a search over a finite grid of sampled schedules and coefficients. Sampled means are estimates at sampled cells; a negative cell is not a proof over the schedules between the cells, and a positive cell is not an optimum. We state that once here and report cell values plainly below.

We score a push-and-unwind manipulator against its zero-impact paired reference (identical seed and schedule, $\lambda = \eta = 0$) over eight seeds, so \texttt{impact\_pnl} attributes profit specifically to having moved the price. \texttt{impact\_boundary\_sweep} sweeps permanent impact (Kyle $\lambda$), temporary impact ($\eta$) and the momentum-follower gain over six grid points per axis; \texttt{size\_response} sweeps the push weight over five points. The evidence commits per-seed impact-PnL vectors at every grid point; intervals below are t-based 95\% CIs over the eight seeds.

Under the precise linear, time-independent permanent-impact assumptions studied by \citet{huberman2004price} and the related no-dynamic-arbitrage conditions of \citet{gatheral2010no}, round-trip price manipulation is ruled out, so the linear sweep is a model-consistency check on the environment rather than a discovery procedure. The three units that follow test what a linear symmetric sweep structurally cannot reach.

\begin{table}[t]
\centering
\caption{Manipulation probe, impact-attributable PnL as fraction of capital per episode, in units of $10^{-4}$, mean over 8 seeds with t-based 95\% CIs. Left: endpoints of each swept axis (6 grid points per axis); every axis is unprofitable everywhere, so no profitability boundary exists to report. Right: the full size-response grid in the push weight.}
\label{tab:f5}
\small
\begin{tabular}{lrrl}
\toprule
Axis & Low end & High end & Boundary \\
\midrule
Kyle $\lambda$ (0 to 0.8) & $-1.3$ & $-31.7\ [-44.3, -19.2]$ & none \\
$\eta$ (0 to 0.8) & $-2.2$ & $-22.4\ [-23.5, -21.4]$ & none \\
Follower gain (0 to 120) & $-2.6$ & $-6.7\ [-10.1, -3.3]$ & none \\
\bottomrule
\end{tabular}
\hspace{1.2em}
\begin{tabular}{lrl}
\toprule
Push wt. & PnL & 95\% CI \\
\midrule
0.10 & $-0.14$ & $[-0.25, -0.03]$ \\
0.25 & $-0.50$ & $[-0.77, -0.24]$ \\
0.50 & $-1.53$ & $[-2.05, -1.00]$ \\
0.75 & $-3.07$ & $[-3.84, -2.30]$ \\
1.00 & $-5.14$ & $[-6.16, -4.12]$ \\
\bottomrule
\end{tabular}
\end{table}

\Cref{tab:f5} and \cref{fig:f5,fig:f5-size} are consistent with the linear model's no-manipulation conditions on the sampled grid. On every axis the pump loses money at every sampled coefficient, with every per-point 95\% CI excluding zero, including the realistic center of the grid ($\lambda = 0.1$, $\eta = 0.05$: impact PnL $-3.4 \times 10^{-4}$, CI $[-4.3, -2.6] \times 10^{-4}$). The size response is bounded and decreasing: the least-bad size is the smallest sampled push weight and the loss grows to $-5.1 \times 10^{-4}$ at full size.

\begin{figure}[tbp]
\centering
\includegraphics[width=0.92\linewidth]{f5-boundaries}
\caption{The manipulation probe, boundary sweeps per impact axis: impact-attributed PnL is negative at every sampled coefficient on every axis.}
\label{fig:f5}
\end{figure}

\begin{figure}[tbp]
\centering
\includegraphics[width=0.92\linewidth]{f5-size-response}
\caption{The manipulation probe, size response in the push weight: the loss is bounded and decreasing in size, so the least-bad push is the smallest one sampled.}
\label{fig:f5-size}
\end{figure}

\subsubsection{Normalized-flow concavity: a falsifiable symmetric probe}\label{sec:f5-concave}

Empirical impact estimates are concave in traded quantity rather than linear \citep{almgren2005direct}, so a probe that can only express linear impact cannot return a positive even in principle. To make it capable of one, the engine supplies an opt-in permanent-impact exponent on dimensionless crowd flow $x=Q/V$:
\[
I_\beta(Q;V)=\lambda\,\operatorname{sign}(x)|x|^\beta,
\qquad x=Q/V.
\]
The temporary own-size term remains $\eta q_i/V$. This normalization is load-bearing: applying a power to raw shares and dividing by $V$ afterwards would change units with $\beta$ and make results depend on the arbitrary share denomination. A regression test rescales capital, flow and $V$ together and requires the nonlinear impact fraction to remain unchanged. At $\beta=1$ the engine uses the original expression $(\lambda Q+\eta q_i)/V$ exactly, preserving every linear golden bit; only $\beta\ne1$ enters the \texttt{powf} path outside the cross-runtime golden guarantee. At fixed $\lambda$ and sub-unit $|x|$, lowering $\beta$ also raises the impact magnitude relative to the linear term, so this is a falsifiability stress test rather than a scale-matched comparison of impact shapes.

We rerun the symmetric push-and-unwind sweep at $\beta\in\{0.5,0.7\}$, eight seeds per point, with paired zero-impact references and pointwise $t$ intervals. No sampled mean is positive. At the canonical calibration, impact-attributed P\&L is $-8.9\times10^{-4}$ (95\% CI $[-19.6,1.8]\times10^{-4}$) for $\beta=0.5$, $-13.5\times10^{-4}$ (CI $[-21.2,-5.9]\times10^{-4}$) for $\beta=0.7$, and $-3.4\times10^{-4}$ (CI $[-4.3,-2.6]\times10^{-4}$) for the linear arm. For $\beta=0.5$, 7 of 23 stored sweep slots have intervals crossing zero; repeated canonical slots reduce these to five distinct configurations. At $\beta=0.7$, 0 of 23 cross zero (\cref{fig:f5-concave}).

\begin{figure}[t]
\centering
\includegraphics[width=\linewidth]{f5-concave}
\caption{Pointwise impact-attributed P\&L intervals for the sampled symmetric schedule under linear and normalized-flow concave impact.}
\label{fig:f5-concave}
\end{figure}

The null identifies the schedule's limitation rather than closing the economic question: profitable constructions under nonlinear permanent impact can be asymmetric, and the next unit searches such shapes.

\subsubsection{Exploratory asymmetric-schedule positive control}\label{sec:f5-positive-control}

Concavity makes slicing matter. For positive dimensionless flow $x=Q/V$, the sum of the $n$ per-slice impact increments before the engine's multiplicative price compounding is
\[
n I_\beta(Q/n;V)=n^{1-\beta} I_\beta(Q;V),
\]
so a slowly accumulated leg and a block liquidation need not cancel when $\beta<1$. The identity is about the impact function, not the compounded terminal price, so it motivates a positive control without guaranteeing that any finite schedule profits after compounding, temporary costs and follower flow.

We cross five up/down duration ratios, three block fractions, three exponents and three market arms: 135 cells, each on eight seeds. All displayed cell intervals are pointwise. Because the reported winner is selected after seeing the grid, its headline inference uses a two-sided Bonferroni familywise 95\% interval over all 135 cells (Student $t$, 7 degrees of freedom); the evidence stores the family size, critical value and adjusted interval.

The best observed cell is the no-follower, zero-temporary-impact arm at $\beta=0.5$: nine bars of uniform accumulation followed by one-bar liquidation yield $+21.8\times10^{-4}$. Its pointwise 95\% interval is $[17.4,26.2]\times10^{-4}$ and its familywise interval remains positive at $[10.0,33.6]\times10^{-4}$. The same sampled schedule with temporary impact yields $+18.4\times10^{-4}$, with familywise interval $[6.4,30.4]\times10^{-4}$. The canonical follower arm has no pointwise-positive cell. Three granularities of the same count are worth separating: of the 135 cells, four have familywise-positive intervals (the $9{:}1$ and $8{:}2$ uniform schedules on each of the two no-follower arms, $+21.8$, $+10.9$, $+18.4$ and $+8.9\times10^{-4}$); they occupy two arms; and the two reported winners are the $9{:}1$ cell on each arm. The fourth of those cells, temporary-impact $8{:}2$, clears zero only at the boundary: its familywise lower bound is $+0.04\times10^{-4}$, so a reader recomputing it at a slightly different critical value would count three rather than four. At $\beta=0.7$ no sampled mean is positive. Among the 45 sampled linear cells, every mean is negative; the $1{:}9$ and $9{:}1$ theory-arm values are numerically near-symmetric rather than bit-identical (\cref{fig:f5-positive-control}).

\begin{figure}[t]
\centering
\includegraphics[width=\linewidth]{f5-positive-control}
\caption{Exploratory asymmetric-schedule search. Shading shows pointwise 95\% intervals; selection-adjusted intervals for the two selected $\beta=0.5$ winners are reported in the text. Four of the 135 cells retain positive familywise intervals, spread over two sampled arms.}
\label{fig:f5-positive-control}
\end{figure}

We confirm the selected theory-arm cell independently on 32 fresh seeds fixed before that confirmation run: mean $+26.3\times10^{-4}$, two-sided 95\% $t$ interval $[23.7,28.9]\times10^{-4}$ (df 31). This is one fixed-cell interval and carries no multiplicity correction for the exploratory search, so the 135-cell Bonferroni interval answers the selected-grid question while the fresh-band interval answers the one-cell replication question. What the positive control establishes is that this instrument can detect a profitable sampled round trip under normalized concave impact. What it leaves open is the shape of the profitable region: the search fixes push size, $\lambda$, $V$, total leg length and most follower gains, and omits short-side, overshooting and longer-hold trips. Conversely, the absence of a positive among 45 sampled linear cells is consistent with the cited linear-impact model under its assumptions.

\subsubsection{Extended sweeps through the positive-control winner}\label{sec:f5-extended-sweeps}

The positive control leaves six axes at their anchor values, so we sweep each of them one at a time through its no-follower, zero-temporary-impact, uniform $9{:}1$ winner: push weight, $\lambda$, total leg length, hold length, short-side trips and follower gain, at $\beta\in\{1,0.5\}$, five values per axis and eight common seeds, for 60 cells. Because the anchor was selected in the earlier 135-cell grid on these same seeds, the familywise intervals are two-sided Bonferroni 95\% $t$ intervals over the global 195-cell family, the antecedent grid plus this extension, rather than over 60 cells alone. Unequal ``uniform'' legs use separate $1/n_{\rm up}$ and $1/n_{\rm down}$ increments, and tests verify that each leg executes exactly one signed peak notional, so changing duration no longer silently changes the amount liquidated.

The results are sharply asymmetric by exponent. Across the 60 extension cells, 23 global-familywise intervals are positive, 31 are negative and 6 cross zero. The longest tested $45{:}5$ theory-arm trip is the largest observed cell, $+135.0\times10^{-4}$ with global-familywise interval $[122.1,148.0]\times10^{-4}$. Interactions between axes, other exponents, other flow scales, transient impact and overshooting round trips remain untested (\cref{fig:f5-extended-sweeps}).

\begin{figure}[t]
\centering
\includegraphics[width=\linewidth]{f5-extended-sweeps}
\caption{Extended manipulation sweeps after per-leg uniform normalization. Shading is the
pointwise 95\% interval; error bars are the global 195-cell Bonferroni familywise 95\%
interval.
The displayed maxima are sampled cells, not optimized values.}
\label{fig:f5-extended-sweeps}
\end{figure}

\subsection{Mechanism realism: adverse-selection markout (F6)}\label{sec:f6}

\texttt{compare\_informed\_vs\_uninformed} runs twenty-four paired episodes in which the identical meta-order flow is sided on the alpha in one leg and on a coin flip in the other, holding schedule, sizes and price path fixed, and reports mean maker markout per filled unit at horizons of 1, 5 and 20 bars. Markout is measured in the tape's absolute price units (the mid starts at 100) per unit of filled quantity; the evidence commits per-episode vectors, and the gap CIs below are t-based 95\% over the twenty-four paired episodes.

\begin{table}[t]
\centering
\caption{Maker markout per filled unit (price units per unit quantity) against informed versus paired-uninformed flow, 24 paired episodes; gap CIs are t-based 95\% over per-episode paired gaps. The gap is the price of trading against information and grows with horizon; the levels stay positive at every horizon, a separate finding discussed in the text.}
\label{tab:f6}
\small
\begin{tabular}{lrrrl}
\toprule
Horizon (bars) & Informed & Uninformed & Gap & Gap 95\% CI \\
\midrule
1 & 0.689 & 0.761 & 0.072 & $[0.052, 0.094]$ \\
5 & 0.489 & 0.791 & 0.302 & $[0.228, 0.383]$ \\
20 & 0.386 & 0.823 & 0.436 & $[0.303, 0.579]$ \\
\bottomrule
\end{tabular}
\end{table}

Two quantities must be separated here: the sign of the \emph{paired estimated gap} and the sign of the \emph{levels}, and they come out on opposite sides.

The paired estimated gap is positive at every horizon (\cref{tab:f6}, \cref{fig:f6}): makers filled by informed flow mark out worse than the paired-uninformed control at $h = 1$, 5 and 20, the per-unit gap widens from 0.072 to 0.436 as the information realizes, and every horizon's paired CI excludes zero, including the smallest ($h = 1$: $[0.052, 0.094]$ over per-episode paired gaps). The loss is in the drift after the fill rather than in the fill itself, which is the signature shape of adverse selection. The claim is scoped to the paired estimated gap, the differential cost of informed over uninformed flow, and not to maker profitability levels.

The level result goes against the environment, and we report it the way we report F4. Makers earn a \emph{positive} mean markout against informed flow at every horizon (0.689, 0.489, 0.386 per unit), which is on the wrong side of the \cref{sec:principles} standard that a realistic market should not let makers profit from informed flow. Three calibration choices produce it: fills are struck at the pre-move mid, so every fill books the full quoted depth as spread capture; the default quote depths are wide relative to the drift the $\alpha = 2.0$ meta-orders realize over 20 bars; and the price path is exogenous, so quotes never widen in equilibrium response to toxicity as they would in \citet{glosten1985bid}. At this calibration the probe therefore measures markout accounting against a fixed drift, in which informed flow costs makers relative to uninformed flow but not in absolute terms. Recalibrating the informed alpha or striking fills at post-impact prices so the level turns negative, and adding a variant where quotes widen with realized toxicity, are the named follow-ups; the first of them is the subsection below.

The single-episode maker decomposition illustrates the exact identity on one seeded episode (seed 0); it is an illustration, not an aggregate finding. The wider quoter (\texttt{maker\_0}, 135 fills) keeps a positive markout at $h = 20$ (0.526 per unit, toxic-fill rate 0.296): the flow is not selecting against it. The tighter quoter (\texttt{maker\_1}, 92 fills) shows the subsidized pattern: aggregate markout is still positive at $h = 20$ (0.197 per unit) but its meta-order fills alone lose 0.197 per unit on 62\% of its filled quantity, with a toxic-fill rate of 0.533 at $h = 20$, so noise flow is paying for the informed flow. Spread capture and adverse drift decompose exactly ($\texttt{spread\_capture} + \texttt{adverse\_drift} = \texttt{markout}$, in price units times quantity): \texttt{maker\_1} captures 82.8 of spread and gives back 64.7 to adverse drift by $h = 20$.

\begin{figure}[t]
\centering
\includegraphics[width=\linewidth]{f6-markouts}
\caption{Markout per filled unit by horizon: informed flow versus the paired-uninformed control.}
\label{fig:f6}
\end{figure}

\subsubsection{The endogenous arm: maker markout when the flow moves the price}\label{sec:f6-endogenous}

The design above fixes the efficient-price path, so it measures the informed-versus-uninformed \emph{gap} but cannot speak to maker profitability \emph{levels}: a maker whose fills never move the price against it books the full quoted depth on every print, and the positive levels reported above are partly that artifact. This arm makes the level question answerable. We rerun the same twenty-four paired episodes with the cleared price endogenous: the taker flow actually filled against the maker ladder moves the mid through the permanent-impact law of the shared-book clearing engine of \cref{sec:env}. Let $S_t$ be the efficient price (the meta-orders' alpha increments plus ordinary volatility, identical across both legs and both arms), $Q_t$ the signed taker quantity filled in bar $t$ (meta-order children plus noise flow, buys positive), $\lambda$ Kyle's permanent-impact coefficient \citep{kyle1985continuous} and $V$ the engine's ADV-like volume normalizer. Makers then quote around
\begin{equation}\label{eq:f6-endo}
\mathrm{mid}_t \;=\; S_t\, M_t,
\qquad
M_{t+1} \;=\; M_t \bigl(1 + \lambda\, Q_t / V\bigr),
\qquad
M_0 = 1,
\end{equation}
the engine's running multiplier, updated after the bar clears so the reference price a maker quotes against is moved only by flow strictly before it. Temporary impact is not a separate term: the maker ladder itself is the temporary component (a taker pays the quoted depth, which is what the engine's Almgren--Chriss $\eta$ term models in the maker-less market \citep{almgren2001optimal}), so the arm adds exactly one mechanism to the exogenous scenario, that a filled unit moves the next reference mid by $\mathrm{mid}_t\,\lambda/V$ in the taker's direction. At $\lambda = 0$ the arm reproduces the exogenous scenario bit for bit, and a cross-check in the test suite feeds the identical flow sequence to the native clearing engine and matches its cleared mids to relative $10^{-12}$, so \cref{eq:f6-endo} is the engine's law and not a paraphrase of it. The calibration sets $V = 240$, one episode's realized taker volume (mean 239 filled units over the 48 exogenous-arm episodes), and the engine's default $\lambda = 0.1$, so a filled unit displaces the mid by $0.042$ price units at the initial mid of 100; a parent order fills a mean 22.4 of its 60 offered units, displacing the mid by a mean $0.81\%$ over its own window (t-based 95\% CI $[0.77, 0.85]\%$ over episodes, sided with the parent in both legs, and zero by construction in the exogenous arm).

\begin{table}[t]
\centering
\caption{Maker markout per filled unit (price units per unit quantity) under the exogenous and endogenous price paths, 24 paired episodes, horizons in bars. Every cell is the mean over per-episode values (episode markout over episode filled quantity, all makers; toxic rate is the share of all maker fills marking out negative), with t-based 95\% CIs, $\mathrm{df}=23$; this is the convention of the committed per-episode vectors and of the gap CIs in \cref{tab:f6}, so the exogenous rows differ from the pooled quantity-weighted aggregates of \cref{tab:f6} in the third decimal. The evidence also serializes paired endogenous-minus-exogenous changes and their CIs. Endogenous arm at $\lambda = 0.1$, $V = 240$.}
\label{tab:f6-endo}
\footnotesize
\setlength{\tabcolsep}{3.5pt}
\begin{tabular}{llrrrrr}
\toprule
Path & $h$ & Informed & Uninformed & Gap & \multicolumn{2}{c}{Toxic-fill rate (inf / uninf)} \\
\midrule
exogenous  & 1  & 0.689 $[0.683, 0.695]$ & 0.762 $[0.740, 0.783]$ & 0.073 $[0.052, 0.094]$ & 0.016 & 0.044 \\
exogenous  & 5  & 0.488 $[0.470, 0.507]$ & 0.794 $[0.720, 0.867]$ & 0.305 $[0.228, 0.383]$ & 0.221 & 0.157 \\
exogenous  & 20 & 0.385 $[0.341, 0.430]$ & 0.827 $[0.711, 0.942]$ & 0.441 $[0.303, 0.579]$ & 0.381 & 0.279 \\
\midrule
endogenous & 1  & 0.622 $[0.614, 0.629]$ & 0.693 $[0.672, 0.715]$ & 0.072 $[0.051, 0.093]$ & 0.034 & 0.056 \\
endogenous & 5  & 0.352 $[0.327, 0.376]$ & 0.652 $[0.579, 0.724]$ & 0.300 $[0.225, 0.375]$ & 0.323 & 0.209 \\
endogenous & 20 & 0.240 $[0.192, 0.289]$ & 0.670 $[0.558, 0.782]$ & 0.430 $[0.294, 0.565]$ & 0.437 & 0.297 \\
\bottomrule
\end{tabular}
\end{table}

Both effects the design predicts appear (\cref{tab:f6-endo}, \cref{fig:f6-endo}). First, the gap moves by at most $0.011$ per unit: $0.072$, $0.300$ and $0.430$ at $h = 1$, 5 and 20 against $0.073$, $0.305$ and $0.441$ in the exogenous arm, every CI excluding zero. Permanent impact is direction-agnostic from the maker's seat, since the maker is on the wrong side of the move its own fill causes whichever way the taker trades, so it cancels out of the informed-minus-uninformed difference; the gap is the information, and the arm confirms it is not an artifact of the fixed path. Second, the levels fall in both legs at every horizon: informed markout drops from $0.385$ to $0.240$ per unit at $h = 20$, uninformed from $0.827$ to $0.670$, and the informed toxic-fill rate rises from $0.381$ to $0.437$. That drop is the impact cost the exogenous design could not book.

\begin{table}[t]
\centering
\caption{Endogenous-arm markout per filled unit at $h = 20$ as the permanent-impact coefficient sweeps, $V = 240$ throughout, the same 24 paired episodes, and exploratory pointwise t-based 95\% CIs ($\mathrm{df}=23$); no sweep-wide multiplicity claim is made. ``Impact/unit'' is the mid displacement per filled unit at the initial mid, $100\,\lambda/V$. The $\lambda = 0$ row is the exogenous arm.}
\label{tab:f6-endo-sweep}
\footnotesize
\setlength{\tabcolsep}{3.5pt}
\begin{tabular}{rrrrrr}
\toprule
$\lambda$ & Impact/unit & Informed & Uninformed & Gap & Toxic (inf) \\
\midrule
0.00 & 0.000 & $+0.385$ $[+0.341, +0.430]$ & $+0.827$ $[+0.711, +0.942]$ & 0.441 $[0.303, 0.579]$ & 0.381 \\
0.05 & 0.021 & $+0.313$ $[+0.267, +0.359]$ & $+0.748$ $[+0.635, +0.862]$ & 0.435 $[0.299, 0.572]$ & 0.415 \\
0.10 & 0.042 & $+0.240$ $[+0.192, +0.289]$ & $+0.670$ $[+0.558, +0.782]$ & 0.430 $[0.294, 0.565]$ & 0.437 \\
0.20 & 0.083 & $+0.095$ $[+0.039, +0.152]$ & $+0.513$ $[+0.403, +0.623]$ & 0.418 $[0.282, 0.554]$ & 0.479 \\
0.30 & 0.125 & $-0.050$ $[-0.117, +0.017]$ & $+0.356$ $[+0.245, +0.467]$ & 0.406 $[0.268, 0.544]$ & 0.507 \\
0.40 & 0.167 & $-0.196$ $[-0.275, -0.117]$ & $+0.198$ $[+0.084, +0.312]$ & 0.393 $[0.251, 0.536]$ & 0.525 \\
\bottomrule
\end{tabular}
\end{table}

The economically meaningful question is whether makers still profit against informed flow once their own fills move the price, and at this calibration the answer is yes: the informed level at $h = 20$ is $0.240$ with pointwise CI $[0.192, 0.289]$, positive but $0.145$ per unit lower than the exogenous arm's. The sweep in \cref{tab:f6-endo-sweep} is exploratory. Its sampled informed level declines from $+0.095$ at $\lambda=0.2$ to $-0.050$ at $\lambda=0.3$ (whose pointwise interval crosses zero) and $-0.196$ at $\lambda=0.4$, while the uninformed level remains positive over these six values and the paired gap changes from $0.441$ to $0.393$. The data therefore bracket an observed sign change between sampled values; they do not identify a crossover at $\lambda\approx0.27$ or attach uncertainty to an interpolated value. The default calibration ($\lambda=0.1$) is one explicitly sampled point below that bracket.

The impact model is the engine's stylized linear permanent impact with a single volume normalizer rather than an estimated impact curve, and linearity is itself a modeling commitment \citep{huberman2004price}; the temporary component is folded into the maker ladder rather than carried as a separate $\eta$ term; fills are still struck at the maker's quote around the pre-move mid; and the quoting side is one follower model, the same fixed Avellaneda--Stoikov-plus-fixed-spread roster as above, which never widens in response to realized toxicity, so this is still not the equilibrium quote response of \citet{glosten1985bid}. What the arm establishes is narrower: on this sampled calibration grid, with the same quotes and information, informed-flow maker profitability declines as flow impact increases, while the paired estimated gap changes modestly. The exogenous and endogenous arms share one code path and one artifact, pinned by a golden hash asserting that the exogenous reports do not move when the endogenous arm runs; the endogenous evidence, per-episode vectors and paired arm changes included, lives under the \texttt{endogenous} key of the same artifact.

\begin{figure}[t]
\centering
\includegraphics[width=\linewidth]{f6-endogenous}
\caption{The endogenous arm of F6. Left: maker markout per filled unit by horizon, informed versus paired-uninformed flow, exogenous versus endogenous path ($\lambda = 0.1$, $V = 240$), with t-based 95\% CIs over 24 paired episodes; paired arm changes are serialized in the evidence. Right: endogenous markout per filled unit at $h = 20$ across an exploratory six-value $\lambda$ sweep, with pointwise 95\% CI bands. Dashed line: the default calibration.}
\label{fig:f6-endo}
\end{figure}

\subsection{Adversarial realism: the in-band inversion attack}\label{sec:predictability}

\Cref{sec:limitations} names a threat the leak-free interface property does not cover: a scenario is a pure function of a public 64-bit seed under a public deterministic transform, so future bars are in principle computable from past bars, and no forbidden operation is needed. This probe converts that named threat into a measurement. Three adversaries predict the next bar's return from a causal prefix only, on the three volatility tiers, 16 seeds per tier (4 symbols, 120 days, 30 warmup bars): an \emph{unconditional baseline} that predicts the expanding prefix-mean return per symbol; an \emph{honest statistical adversary}, a ridge AR(5) refit on the prefix at every bar, with no knowledge of the generator; and a \emph{known-seed oracle}, which replays the generator from the true seed and therefore predicts every future bar exactly, the upper bound on any inversion attack. We also run each predictor as a frictionless unit-gross long/short sign-following policy scored through the public \texttt{score\_run} kernel, and report the kernel's deflated-Sharpe probability.

\paragraph{The honest adversary extracts almost nothing.} All ``$\pm$'' quantities here are across-seed standard deviations over the 16 seeds, computed with the population convention (\texttt{ddof} 0), not confidence intervals. On Calm the ridge AR reaches $55.4\% \pm 2.8\%$ directional accuracy against the baseline's $54.1\% \pm 2.4\%$: it detects the tier's designed momentum autocorrelation, worth about one descriptive accuracy point over drift-following. On Hard and Extreme the AR accuracies are $51.2\% \pm 3.3\%$ and $50.6\% \pm 2.1\%$, with MSE no better than the unconditional baseline (Hard: $1.24\times10^{-4}$ vs.\ $1.22\times10^{-4}$; Extreme: $1.27\times10^{-3}$ vs.\ $1.17\times10^{-3}$). No hypothesis test or equivalence margin was preregistered, so these summaries support only the descriptive statement that the adversary lands near the unconditional baseline. One post-hoc check is worth reporting rather than hiding: a paired $t$-test of the ridge AR against the baseline over the same 16 Calm seeds gives $t(15) = 3.36$, $p = 0.004$, so the one-point Calm edge over drift-following is larger than seed noise, exactly as the tier's disclosed momentum structure predicts; the same test is inconclusive on Hard ($p = 0.08$) and Extreme ($p = 0.40$).

\paragraph{The oracle bound is total, and a bounded seed search reaches it.} The known-seed oracle scores $100\%$ accuracy, zero MSE, and DSR $1.00$ on every tier: determinism means the gap between an honest adversary and a successful inverter is the entire distance from coin flipping to omniscience, so the question is only the cost of inversion. We implement the strongest cheap version, a brute scan over a bounded seed band of width $2^{16}$ centered on the published \texttt{start\_level}, matching each candidate's first-bar closes against the single observed opening bar. The scan recovers the true seed in 16 of 16 trials with zero collisions, from a prefix of one bar, at $14.7\,\mu$s per candidate ($\approx 0.96$\,s per scan); every recovered seed then reproduces the full 30-bar observed prefix on its tier. When a deployment uses a bounded public \texttt{(start\_level, num\_levels)} specification, this in-band inversion is a table scan rather than cryptanalysis, and its scenarios should be treated as known to a motivated adversary (\cref{fig:predictability}). Extrapolating this Python harness's measured per-candidate cost to the unbounded $2^{64}$ space gives $\sim 8.6$ million CPU-years. The figure is harness-specific: a native scanner would be orders of magnitude cheaper per candidate, which moves the number but not the conclusion that brute-force inversion of an unpublished 64-bit seed is infeasible by enumeration.

\begin{figure}[t]
\centering
\includegraphics[width=\linewidth]{predictability}
\caption{Next-bar directional accuracy by tier. The causal AR adversary remains near the unconditional baseline; knowing or recovering the seed gives perfect prediction.}
\label{fig:predictability}
\end{figure}

\paragraph{What remains open.} The probe brackets the attack but does not close it. Between the honest regression (no edge beyond disclosed structure) and the band scan (total recovery given a bounded public band) lies the case that matters for unbounded deployments: analytic inversion of SplitMix64 \citep{steele2014splittable} through the mul/add/div/max price transform from observed reals. We implement and test the narrow algebraic step: SplitMix64's 64-bit finalizer is exactly invertible when its complete output word is known. The generator's \texttt{next\_unit}, however, exposes only the upper 53 bits as a float, leaving $2^{11}$ candidate post-increment states from one exact unit observation before any price arithmetic is considered. The public interface reveals neither that unit nor an unrounded finalizer output: it reveals prices after data-dependent arithmetic, floors and serialization. Deciding among those candidates from observed prices remains the untested cryptanalytic step; the new primitive reduces the question without claiming a tape-level break. This probe therefore supports only the narrower conclusion that resistance to this attack rests on seed secrecy and entropy rather than transform hardness, and that published bounded seed bands forfeit both.

\subsubsection{Sealed evaluation seeds: defeating the bounded-band scanner}\label{sec:sealed-seeds}

The public held-out set is disjoint from training but enumerable, so the band scanner above recovers its scenario identities without violating the observation interface. The opt-in mitigation derives each named slot from a secret salt $k$:
\[
\mathrm{seed}_i=\texttt{EVAL\_SEED\_BASE}+
\bigl(F(k,i)\bmod(2^{64}-1-\texttt{EVAL\_SEED\_BASE})\bigr).
\]
The shipped $F$ combines FNV-1a/64 with public SplitMix64 finalizers. It is deterministic and pinned across Rust and Python, and the tested salts and slots produce held-out-band seeds distinct from the public set. No universal injectivity claim is possible for a finite 64-bit mapping, and this construction is PRF-style rather than a certified MAC.

The intended workflow is commit--reveal, as \cref{sec:protocol-custody} prescribes it. Before evaluation, the host publishes SHA-256$(k)$ and withholds $k$; after scoring, it reveals $k$ so anyone can recompute the named seeds and replay the run. SharpeBench's forward-window format can bind that hash before entries arrive. This paper's evidence is a simulated workflow containing the revealed salt, not proof that any live evaluation used secret custody, a non-reuse policy or a reveal deadline.

On 16 tested slots, the scanner over the public $2^{16}$ band recovers 16/16 public seeds and 0/16 sealed seeds while the salt is withheld; revealing the salt replays 16/16. Under a uniform-seed approximation, that scanner covers about $3.6\times10^{-15}$ of the derived band, so its expected number of recoveries over 16 slots is $5.7\times10^{-14}$: the 0/16 outcome is what the coverage arithmetic already implies, and the run verifies the implementation rather than testing an uncertain hypothesis. The result is therefore narrow: high-entropy salt secrecy defeats this particular bounded-band enumeration budget. It is not a cryptographic security proof, and it does not answer tape-level inversion after observations begin.
